\exercisesection{4}

\begin{enumerate}
\def\labelenumi{\arabic{enumi}.}
\tightlist
\item
  设 A 为整闭 Noether 整环，V 为 A 的分式域上的有限秩向量空间。证明：若 \((\mathbf{M}_{\lambda})\) 是 V 中都包含同一个格 N 的自反格的任一族，则格 \(M = \bigcap_{\lambda} M_{\lambda}\) 是自反的（考虑对偶格 \(M_{\lambda}^{*}\)）。
\end{enumerate}

* 2. 设 A 为整闭 Noether 整环，E、F 为两个有限生成 A-模。假设 E 是无扭的，\(\mathrm{Hom}_{\mathrm{A}}(\mathrm{E},\mathrm{F})\) 是自反的，且 \(\mathrm{Ext}_{\mathrm{A}}^{1}(\mathrm{E},\mathrm{F}) = 0\)。证明 F 是自反的。（先证明 F 是无扭的；若 \(\mathrm{T} = \mathrm{F}^{**}/c_{\mathrm{F}}(\mathrm{F})\)，利用正合列以两种方式计算 \(\mathrm{Ass}(\mathrm{Hom}_{\mathrm{A}}(\mathrm{E},\mathrm{T}))\)：

\[
0 \rightarrow \operatorname{Hom} (\mathrm{E}, \mathrm{F}) \rightarrow \operatorname{Hom} (\mathrm{E}, \mathrm{F} ^ {* *}) \rightarrow \operatorname{Hom} (\mathrm{E}, \mathrm{T}) \rightarrow 0
\]

以及第 IV 章 § 1 第 4 节的命题 10。） *

\begin{enumerate}
\def\labelenumi{\arabic{enumi}.}
\setcounter{enumi}{2}
\item
  设 A 为整闭 Noether 整环，K 为其分式域，M 为 K 上某有限秩向量空间的自反格，L 为包含 M 的 V 的自由格。证明存在 V 的自由格 \(L_{1}\) 使 \(M = L \cap L_{1}\)。（考虑使 \(L_{p} \neq M_{p}\) 的 A 的高度为 1 的素理想 p 所成的有限集 I，以及主理想整环 \(S^{-1}A\)，其中 \(S = \bigcap_{p \in I} (A - p)\)；证明存在 V 的自由格 \(L_{0}\)，使 \(M_{p} = (L_{0})_{p}\) 对所有 \(p \in I\) 成立，且存在 \(s \in S\)，使 \(M \subset s^{-1}L_{0} \subset L_{1}\)。）
\item
  设 k 为域，A 为多项式环 \(k[X, Y]\)。
\end{enumerate}

\begin{enumerate}
\def\labelenumi{(\alph{enumi})}
\item
  设 \((e_1, e_2)\) 为 A-模 \(A^2\) 的典范基，E 为 \(A^2\) 的由 \((X - Y)e_1\)、\(e_1 + Xe_2\) 与 \(e_1 + Ye_2\) 生成的 A-子模；设 F 为 E 的由 \((X - Y)^2e_1\) 生成的单生成子模。证明 A-模 \(M = E / F\) 不是其扭子模与一个无扭模的直和。
\item
  证明扭 A-模 A/AXY 不是形如 \(A/P_{i}^{n_{i}}\) 的单生成子模的直和，其中诸 \(P_{i}\) 是高度为 1 的素理想。
\end{enumerate}

\begin{enumerate}
\def\labelenumi{\arabic{enumi}.}
\setcounter{enumi}{4}
\tightlist
\item
  设 A 为整闭 Noether 整环，M 为有限生成 A-模。证明：若 a、b 是 A 的两个理想，则 A-模
\end{enumerate}

\[
((a \mathbf {M}) \cap (\mathfrak {b} \mathbf {M})) / (a \cap \mathfrak {b}) \mathbf {M}
\]

是伪零的；给出一个使它不为零的例子。

\begin{enumerate}
\def\labelenumi{\arabic{enumi}.}
\setcounter{enumi}{5}
\tightlist
\item
  设 \(k\) 为域，\(\mathbf{B}\) 为多项式环 \(k[\mathbf{X},\mathbf{Y}]\)，\(\mathbf{A}\) 为子模 \(k[\mathbf{X}^2,\mathbf{XY},\mathbf{Y}^2]\)（包含于 \(\mathbf{B}\)）。
\end{enumerate}

\begin{enumerate}
\def\labelenumi{(\alph{enumi})}
\item
  证明 A 是整闭 Noether 整环，且 B 是有限生成 A-模。
\item
  证明 A 的理想 \(p = AX^{2} + AXY\) 是高度为 1 的素理想（从而是除性的），但 B-模 \(p \otimes_{A} B\) 是非零 \(\neq 0\) 的扭模，从而不是自反的；B 的理想 pB 不是除性的，典范映射 \(p \otimes_{A} B \to pB\) 不是单射，且 p 不是平坦 A-模。
\end{enumerate}

\begin{enumerate}
\def\labelenumi{\arabic{enumi}.}
\setcounter{enumi}{6}
\tightlist
\item
  设 A 为整闭 Noether 整环，E 为有限生成无扭 A-模，\(\mathbf{E}^*\) 为其对偶模。
\end{enumerate}

\begin{enumerate}
\def\labelenumi{(\alph{enumi})}
\item
  证明典范同态 \(\mathbf{E}^{*}\otimes_{\mathbf{A}}\mathbf{E}\to \mathrm{End}_{\mathbf{A}}(\mathbf{E})\) 是伪同构。
\item
  由 (a) 推出：E 为投射 A-模的充分必要条件是 \(E^{*} \otimes_{A} E\) 为自反 A-模。（注意，若 \(E^{*} \otimes_{A} E\) 自反，则典范同态 \(E^{*} \otimes_{A} E \to \operatorname{End}_{A}(E)\) 是双射。）
\end{enumerate}

¶* 8. 设 A 为整闭 Noether 整环，\(\mathbf{M}_1, \mathbf{M}_2\) 为两个有限生成 A-模。

\begin{enumerate}
\def\labelenumi{(\alph{enumi})}
\item
  证明 A-模 \(\operatorname{Tor}_i^{\mathbf{A}}(\mathbf{M}_1, \mathbf{M}_2)\)、\(\operatorname{Ext}_{\mathbf{A}}^{i}(\mathbf{M}_1, \mathbf{M}_2)\) 对 \(i \geqslant 2\) 都是伪零的（归结为 \(\mathbf{A}\) 是主理想整环的情形）。
\item
  若 \(\mathbf{M}_1\) 是无扭的，证明 \(\operatorname{Tor}_1^{\mathrm{A}}(\mathbf{M}_1, \mathbf{M}_2)\) 与 \(\operatorname{Ext}_A^1(\mathbf{M}_1, \mathbf{M}_2)\) 是伪零的（同样的方法）。
\item
  若 \(\mathbf{M}_1\) 是扭 A-模，证明（用第 5 节的记号）：
\end{enumerate}

\[
\chi (\mathbf {M} _ {1} \otimes_ {\mathbf {A}} \mathbf {M} _ {2}) - \chi (\operatorname{Tor} _ {1} ^ {\mathbf {A}} (\mathbf {M} _ {1}, \mathbf {M} _ {2})) = r (\mathbf {M} _ {2}) \chi (\mathbf {M} _ {1})
\]

\[
\chi (\operatorname{Hom} _ {\mathbf {A}} (\mathbf {M} _ {1}, \mathbf {M} _ {2})) - \chi (\operatorname{Ext} _ {\mathbf {A}} ^ {1} (\mathbf {M} _ {1}, \mathbf {M} _ {2})) = - r (\mathbf {M} _ {2}) \chi (\mathbf {M} _ {1})
\]

\[
\chi (\mathrm{Hom} _ {\mathrm{A}} (\mathrm{M} _ {2}, \mathrm{M} _ {1})) - \chi (\mathrm{Ext} _ {\mathrm{A}} ^ {1} (\mathrm{M} _ {2}, \mathrm{M} _ {1})) = r (\mathrm{M} _ {2}) \chi (\mathrm{M} _ {1}).
\]

\begin{enumerate}
\def\labelenumi{(\alph{enumi})}
\setcounter{enumi}{3}
\tightlist
\item
  证明对有限生成 A-模的每个有序对 \(M_{1}\)、\(M_{2}\)：
\end{enumerate}

\[
\begin{array}{r} c (\mathbf {M _ {1}} \otimes_ {\mathbf {A}} \mathbf {M _ {2}}) - c (\mathrm{Tor} _ {\mathbf {1}} ^ {\mathbf {A}} (\mathbf {M _ {1}}, \mathbf {M _ {2}})) = r (\mathbf {M _ {1}}) c (\mathbf {M _ {2}}) + r (\mathbf {M _ {2}}) c (\mathbf {M _ {1}}) \\ c (\mathrm{Hom} _ {\mathbf {A}} (\mathbf {M _ {1}}, \mathbf {M _ {2}})) - c (\mathrm{Ext} _ {\mathbf {A}} ^ {1} (\mathbf {M _ {1}}, \mathbf {M _ {2}})) = r (\mathbf {M _ {1}}) c (\mathbf {M _ {2}}) - r (\mathbf {M _ {2}}) c (\mathbf {M _ {1}}). \end{array}
\]

\begin{enumerate}
\def\labelenumi{\arabic{enumi}.}
\setcounter{enumi}{8}
\tightlist
\item
  设 k 为域，A 为多项式环 \(k[X,Y]\)。在 A-模 \(M = A^{2}\) 上考虑满足 \(f(e_{1}) = X\)、\(f(e_{2}) = Y((e_{1}, e_{2})\) 为典范基）的线性形式 f。证明 f 的核 L 是单生成自由 A-模，但商模 M/L（它同构于 A 的一个理想）不是自反的。
\end{enumerate}

¶ 10. 设 A 为交换环。对有限生成自由 A-模 \(L = A^n\) 的每个子模 R，用 \(c_1(R)\) 表示由诸 \(\langle x, x^* \rangle\) 生成的理想，其中 \(x\) 取遍 R，\(x^*\) 取遍对偶 \(L^*\)；记 \(c_k(R) = c_1\left(\operatorname{Im}\left(\bigwedge^k R\right)\right)\)，其中 \(\operatorname{Im}\left(\bigwedge^k R\right)\) 是 R 的 \(k\) 次外幂在 \(\bigwedge^k L\) 中的典范像。若 \(R_1 \subset R_2\) 是 L 的两个子模，则 \(c_k(R_1) \subset c_k(R_2)\) 对所有 \(k\) 成立。

\begin{enumerate}
\def\labelenumi{(\alph{enumi})}
\item
  设 M 为有限生成 A-模；M 同构于一个商模 L/R，其中对某个适当的 n 有 \(L = A^{n}\)。证明：理想序列 \((\mathfrak{a}_{k})_{k \geqslant 0}\)（满足 \(\mathfrak{a}_{k} = \mathfrak{c}_{n-k}(R)\)，这里 \(k \leqslant n\)；\(\mathfrak{a}_{k} = A\)，这里 k \textgreater{} n）与 M 表成 L/R 的方式无关。（先考虑 L/R 与 L/R' 同构的情形；然后注意，\(A^{n}/R\) 同构于 \(A^{n+h}/(R \times A^{h})\)，对所有 h \textgreater{} 0。）我们记 \(\mathfrak{d}_{k}(M) = \mathfrak{a}_{k}\)（对所有 \(k \geqslant 0\)），并称它们为与 M 相伴的行列式理想。于是 \(\mathfrak{d}_{k}(M) \subset \mathfrak{d}_{k+1}(M)\) 对 \(k \geqslant 0\) 成立。
\item
  若 \(\mathbf{A}\) 是主理想整环且 \(\mathbf{M} = \mathbf{L} / \mathbf{R}\)，其中 \(\mathbf{L}\) 是自由且有限生成的，证明理想 \(\mathfrak{c}_{k + 1}(\mathbb{R})(\mathfrak{c}_k(\mathbb{R}))^{-1}\) 是 \(\mathbf{R}\) 在 \(\mathbf{L}\) 中的不变因子。
\item
  设 \(r = \gamma(\mathbf{M})\) 为 \(\mathbf{M}\) 的生成系的最小基数，\(r_0\) 为最小整数 \(h\)，使 \(\mathfrak{d}_h(\mathbf{M}) = \mathbf{A}\)。证明 \(r_0 \leqslant r\)，并给出一个使 \(r_0 < r\) 的例子（取 \(\mathbf{A}\) 为 Dedekind 整环）。
\item
  若 \(a\) 是 \(\mathbf{M}\) 的零化子，证明 \(\mathfrak{d}_0(\mathbf{M}) \subset \mathfrak{a}\)，且 \(\mathfrak{a}^{r - k} \subset \mathfrak{d}_k(\mathbf{M})\) 对 \(k \leqslant r = \gamma (\mathbf{M})\) 成立。
\item
  设 \(\mathbf{M}\) 是子模 \(\mathbf{M}_i\)（\(1 \leqslant i \leqslant h\)）的直和。证明
\end{enumerate}

\[
\mathfrak {d} _ {k} (\mathbf {M}) = \sum \mathfrak {d} _ {k _ {1}} (\mathbf {M} _ {1}) \dots \mathfrak {d} _ {k _ {n}} (\mathbf {M} _ {n})
\]

其中和式取遍有限序列 \((k_{i})_{1\leqslant i\leqslant h}\)（满足 \(\sum_{i = 1}^{h}k_{i} = k.\)）

\begin{enumerate}
\def\labelenumi{(\alph{enumi})}
\setcounter{enumi}{5}
\tightlist
\item
  若 N 是 M 的有限生成子模，则 \(\mathfrak{d}_{k}(\mathbf{M}/\mathbf{N}) \subset \mathfrak{d}_{k}(\mathbf{M})\) 对所有 \(k \geqslant 0\) 成立。证明：
\end{enumerate}

\[
\sum_ {j = 0} ^ {k} \mathfrak {d} _ {j} (\mathrm{N}) \mathfrak {d} _ {k - j} (\mathrm{M} / \mathrm{N}) \subset \mathfrak {d} _ {k} (\mathrm{M}).
\]

\begin{enumerate}
\def\labelenumi{(\alph{enumi})}
\setcounter{enumi}{6}
\tightlist
\item
  设 K 为域，A 为多项式环 K{[}X, Y, Z{]}，m 为极大理想 AX + AY + AZ，q 为由 X、Y²、YZ 与 Z² 生成的理想。考虑 A-模 M = A/q 及其子模 N = m/q。证明 \(\mathfrak{d}_{0}(M) = \mathfrak{q}, \mathfrak{d}_{1}(M) = A, \mathfrak{d}_{0}(N) = m^{2}\) 且 \(\mathfrak{d}_{1}(N) = m\)。
\end{enumerate}

\begin{enumerate}
\def\labelenumi{\arabic{enumi}.}
\setcounter{enumi}{10}
\tightlist
\item
  设 A 为 Krull 整环，M、N 为 A 的分式域上有限秩 n 向量空间 V 中的两个有限生成格。对 A 中每个 \(c \neq 0\)（设它满足 \(cN \subset M\)），考虑行列式理想 \(\mathfrak{d}_{k}(M/cN)\)（习题 10）并记：
\end{enumerate}

\[
d _ {k} (\mathbf {N}, \mathbf {M}) = \operatorname{div} \left(\mathfrak {d} _ {k} (\mathbf {M} / c \mathbf {N})\right) - (n - k) \operatorname{div} (c).
\]

\begin{enumerate}
\def\labelenumi{(\alph{enumi})}
\item
  证明 \(d_{k}(\mathbf{N}, \mathbf{M})\) 不依赖于使 \(cN \subset M\) 的 c 的选取；\(d_{k}(\mathbf{N}, \mathbf{M})\) 称为 N 关于 M 的指标为 k 的行列式因子。
\item
  \(d_{k}(\mathbf{N},\mathbf{M})\geqslant d_{k + 1}(\mathbf{N},\mathbf{M})\) 且 \(d_{n}(\mathbf{N},\mathbf{M}) = 0\)。证明：若记
\end{enumerate}

\[
e _ {k} (\mathrm{N}, \mathrm{M}) = d _ {n - k} (\mathrm{N}, \mathrm{M}) - d _ {n - k + 1} (\mathrm{N}, \mathrm{M}),
\]

则

\[
e _ {k} (\mathbf {N}, \mathbf {M}) \leqslant e _ {k + 1} (\mathbf {N}, \mathbf {M})
\]

对 \(1 \leqslant k \leqslant n\) 成立；除子 \(e_k(\mathbf{N}, \mathbf{M})\) 称为 \(\mathbf{N}\) 关于 \(\mathbf{M}\) 的不变因子（归结为 \(\mathbf{A}\) 是主理想整环的情形）。

\begin{enumerate}
\def\labelenumi{(\alph{enumi})}
\setcounter{enumi}{2}
\item
  证明 \(d_0(\mathbf{N}, \mathbf{M}) = \chi(\mathbf{M}, \mathbf{N})\)（第 5 节）。
\item
  若 M 是 V 中的格而 N 是 V 的秩为 q \textless{} n 的子空间 W 中的格，则 N 关于 \(M \cap W\) 的不变因子称为 N 关于 M 的不变因子。说明《代数》第 VII 章 § 4 习题 8 至 10 与 14 至 16 可以怎样推广（在某些情形必要时假设 A 是整闭 Noether 整环）。
\end{enumerate}

\begin{enumerate}
\def\labelenumi{\arabic{enumi}.}
\setcounter{enumi}{11}
\tightlist
\item
  设 A 为整闭 Noether 整环，M、N 为两个同秩 r 的无扭有限生成 A-模，且 \(N \subset M\)；设 \(j: N \to M\) 为典范单射。
\end{enumerate}

\begin{enumerate}
\def\labelenumi{(\alph{enumi})}
\tightlist
\item
  证明对所有 \(k\)，\(\bigwedge^{k}j\colon\bigwedge^{k}\mathrm{N}\to\bigwedge^{k}\mathrm{M}\) 是伪单射的，\(\operatorname{Coker}\left(\bigwedge^{k}j\right)\) 是扭 A-模，且：
\end{enumerate}

\[
\chi \left(\operatorname{Coker} \left(\bigwedge^ {k} j\right)\right) = \binom {r - 1} {k - 1} \chi (\operatorname{Coker} j).
\]

\begin{enumerate}
\def\labelenumi{(\alph{enumi})}
\setcounter{enumi}{1}
\tightlist
\item
  利用 (a) 证明：若 M 是无扭有限生成 A-模，则 \(\bigwedge^{k}\mathbf{M}\) 的扭子模对所有 \(k\) 都是伪零的，且
\end{enumerate}

\[
c \left(\bigwedge^ {k} \mathbf {M}\right) = \binom {r - 1} {k - 1} c (\mathbf {M}),
\]

其中 \(\mathbf{M}\) 的秩为 \(r\)；特别地，\(c\left(\bigwedge^{r}\mathbf{M}\right) = c(\mathbf{M})\)。

\begin{enumerate}
\def\labelenumi{\arabic{enumi}.}
\setcounter{enumi}{12}
\tightlist
\item
  设 \(k\) 为域，\(A\) 为多项式环 \(k[X, Y]\)。
\end{enumerate}

\begin{enumerate}
\def\labelenumi{(\alph{enumi})}
\item
  设 \(m\) 为极大理想 \(AX + AY\)（在 \(A\) 中）；证明存在从 \(m\) 到 \(A\) 的伪同构，但不存在从 \(A\) 到 \(m\) 的伪同构。
\item
  设 \(\mathfrak{m}'\) 为极大理想 \(\mathbf{A}(\mathbf{X} - 1) + \mathbf{A}\mathbf{Y}\)（在 \(\mathbf{A}\) 中）；证明 \(c(\mathfrak{m}) = c(\mathfrak{m}') = 0\)，但既不存在从 \(\mathfrak{m}\) 到 \(\mathfrak{m}'\) 的伪同构，也不存在从 \(\mathfrak{m}'\) 到 \(\mathfrak{m}\) 的伪同构。
\item
  设 \(\mathfrak{p} = \mathbf{AX}\)，\(\mathfrak{q} = \mathbf{AY}\)，它们是高度为 1 的素理想。在 A-模 \(\mathbf{L} = \mathbf{A}^2\) 中考虑子模 \(\mathbf{M} = \mathfrak{pe}_1 \oplus \mathfrak{qe}_2\) 与 \(\mathbf{N} = \mathbf{Ae}_1 \oplus \mathfrak{pqe}_2\)（\(e_1, e_2\) 为 \(\mathbf{L}\) 的典范基的向量）。证明 \(\mathbf{M}\) 与 \(\mathbf{N}\) 同构，且存在从 \(\mathbf{L} / \mathbf{M}\) 到 \(\mathbf{L} / \mathbf{N}\) 的伪同构以及从 \(\mathbf{L} / \mathbf{N}\) 到 \(\mathbf{L} / \mathbf{M}\) 的伪同构，但不存在从 \(\mathbf{L}\) 到自身、把 \(\mathbf{M}\) 映到 \(\mathbf{N}\) 或把 \(\mathbf{N}\) 映到 \(\mathbf{M}\) 的伪同构（注意，借助命题 10，从 \(\mathbf{L}\) 到自身的伪同构必然是 \(\mathbf{L}\) 的自同构）。
\end{enumerate}

¶ 14. 设 A 为 Prüfer 整环（§ 2 习题 12），M、N 为 A 的分式域上有限秩 n 向量空间 V 中的两个有限生成格。对 A 中每个 \(c \neq 0\)（设它满足 \(cN \subset M\)），考虑行列式理想 \(\mathfrak{d}_{k}(M/cN)\)（习题 10）并记：

\[
\mathfrak {d} _ {k} (\mathbf {N}, \mathbf {M}) = c ^ {k - n} \mathfrak {d} _ {k} (\mathbf {M} / c \mathbf {N}).
\]

\begin{enumerate}
\def\labelenumi{(\alph{enumi})}
\item
  证明 \(\mathfrak{d}_k(\mathbf{N},\mathbf{M})\) 不依赖于 \(c\) 的选取，只要 \(c\mathbf{N}\subset \mathbf{M}\)；它们称为 \(\mathbf{N}\) 关于 \(\mathbf{M}\) 的行列式理想。
\item
  \(\mathfrak{d}_k(\mathbf{N},\mathbf{M})\subset \mathfrak{d}_{k + 1}(\mathbf{N},\mathbf{M})\) 且 \(\mathfrak{d}_n(\mathbf{N},\mathbf{M}) = \mathbf{A}\)。证明：若记 \(\mathfrak{e}_k(\mathbf{N},\mathbf{M}) = \mathfrak{d}_{n - k}(\mathbf{N},\mathbf{M})(\mathfrak{d}_{n - k + 1}(\mathbf{N},\mathbf{M}))^{-1}\)，则整理想 \(\mathfrak{e}_k(\mathbf{N},\mathbf{M})\) 满足 \(\mathfrak{e}_k(\mathbf{N},\mathbf{M})\supset \mathfrak{e}_{k + 1}(\mathbf{N},\mathbf{M})\)（对 \(1\leqslant k\leqslant n\)）；有限生成理想 \(\mathfrak{e}_k(\mathbf{N},\mathbf{M})\) 称为 \(\mathbf{N}\) 关于 \(\mathbf{M}\) 的不变因子（考虑 \(\mathrm{A_m}\) -模 \(\mathrm{M_m}\) 与 \(\mathrm{N_m}\)，其中 \(\mathfrak{m}\) 取遍 \(\mathrm{A}\) 的每个极大理想）。
\item
  若 M 是 V 的格而 N 是 V 的秩为 q \textless{} n 的子空间 W 中的格，则 N 关于 M ∩ W 的不变因子称为 N 关于 M 的不变因子。说明《代数》第 VII 章 § 4 习题 8 至 10 与 14 至 16 可以怎样推广。
\end{enumerate}

\begin{enumerate}
\def\labelenumi{\arabic{enumi}.}
\setcounter{enumi}{14}
\tightlist
\item
  设 A 为多项式环 Z{[}X, Y, T, U, V, W{]}，L 为自由 A-模 \(A^{4}\)，\((e_{i})_{1 \leq i \leq 4}\) 为其典范基，M 为 L 的由四个向量 \(Xe_{1}\)、\(Ye_{2}\)、\(Te_{3} + Ue_{4}\)、\(Ve_{3} + We_{4}\) 生成的子模。证明
\end{enumerate}

\[
\mathfrak {d} _ {1} (\mathrm{L} / \mathrm{M}) \mathfrak {d} _ {3} (\mathrm{L} / \mathrm{M}) \notin (\mathfrak {d} _ {2} (\mathrm{L} / \mathrm{M})) ^ {2}
\]

（与习题 14 (b) 比较）。

¶ 16. (a) 设 A 为 Prüfer 整环，M 为 A 的分式域 K 上有限秩 n 向量空间 V 中的有限生成格。证明存在 V 的基 \((e_{i})_{1 \leqslant i \leqslant n}\) 与 n 个分式理想 \(a_{i}\)（\(1 \leqslant i \leqslant n\)），使 M 等于诸 \(a_{i}e_{i}\) 的直和。（对 n 作归纳论证：若 \((u_{i})_{1\leq i\leq n}\) 是 V 的一组基，考虑有限生成分式理想 \(b_{1}\)（由 M 中元素的 \(u_{1}\) -坐标生成）；通过考虑 A-模 \(b_{1}^{-1}M\)，归结为 \(b_{1}=A\) 的情形。）

由此推出，若 W 是 V 的子空间，则 M ∩ W 是 M 的直和项。

\begin{enumerate}
\def\labelenumi{(\alph{enumi})}
\setcounter{enumi}{1}
\tightlist
\item
  设 M 是 \(A^{n}\) 的子模，等于诸 \(a_{i}e_{i}\) 的直和，其中诸 \(a_{i}\) 是 A 的有限生成（整）理想，且 \((e_{i})_{1\leqslant i\leqslant n}\) 是 \(A^{n}\) 的典范基。证明存在基 \((u_{i})_{1\leqslant i\leqslant n}\)（在 \(A^{n}\) 中）与 A 的理想 \(b_{i}\)，使 M 是诸 \(b_{i}u_{i}\) 的直和，且 \(b_{i}\) 整除 \(b_{i+1}\)（对 \(1\leqslant i\leqslant n-1\)）。（归结为 n=2 且 \(a_{1}+a_{2}=A\) 的情形。）
\end{enumerate}

\begin{enumerate}
\def\labelenumi{\arabic{enumi}.}
\setcounter{enumi}{16}
\tightlist
\item
  设 k 为域，A 为多项式环 \(k[X, Y]\)，L 为自由 A-模 \(A^{2}\)，\((e_{1}, e_{2})\) 为 L 的典范基。设 M 为 L 的由两个向量 \((X + 1)e_{1} + Ye_{2}, Ye_{1} + Xe_{2}\) 生成的子模。证明不存在从 L 到自身的伪同构，它限制在 M 上是从 M 到形如 \(au + bv\) 的子模 N 的伪同构（其中 \((u, v)\) 是 A-模 L 的基，a、b 是 A 的两个理想），或者限制在 N 上是从 N 到 M 的伪同构。（通过考虑行列式理想（习题 10）并注意 M 是自反的，可以只限于 N 也自反的情形，而此时必然有 a = A、b = AP，其中
\end{enumerate}

\[
\mathrm{P} (\mathrm{X}, \mathrm{Y}) = \mathrm{X} (\mathrm{X} + 1) - \mathrm{Y} ^ {2}
\]

是 A 的一个极端元；最后证明，对 L 的每组基 \((u, v)\)，M ∩ Au 既不能包含 Au，也不能包含于 APu。）

¶ 18. 设 A 为 Dedekind 整环，K 为其分式域，M、N 为 K 上有限秩 n 向量空间 V 中的两个格。设 \(\mathfrak{e}_{k} = \mathfrak{e}_{k}(N, M)\) 为 N 关于 M 的不变因子（习题 14）。证明存在 V 的基 \((u_{i})_{1 \leqslant i \leqslant n}\)，使 M 等于直和 \(\bigoplus_{i} a_{i} u_{i}\)，其中诸 \(a_{i}\) 是分式理想，而 N 等于直和 \(\bigoplus_{i} e_{i} a_{i} u_{i}\)。（利用主理想整环上有限生成模的理论（《代数》第 VII 章 § 4 第 2 节定理 1）以及单模群 \(\mathbf{SL}(n, \mathbf{A})\) 中的逼近定理（§ 2 第 4 节）。）

把这一结果推广到 A 是一族右有向的 Dedekind 整环子环之并的情形（例如 Dedekind 整环在其分式域的代数闭包中的整闭包）。

\begin{enumerate}
\def\labelenumi{\arabic{enumi}.}
\setcounter{enumi}{18}
\item
  设 A 为 Dedekind 整环，a、b 为 A 的分式理想。证明 A-模 A ⊕ ab 与 a ⊕ b 同构。
\item
  设 A 为 Dedekind 整环，P 为投射 A-模，N 为 P 的子模。证明 N 是投射的，且是同构于 A 的理想的模的直和。（可以只限于 P 是自由模的情形。然后照习题 16 那样进行，使用以《代数》第 VII 章 § 3 定理 1 为基础的超限归纳论证。）
\item
  设 P 为 Dedekind 整环 A 上的投射模。证明若 P 不是有限生成的，则它是自由的。（据习题 20，可把 P 写成无限直和 \(\bigoplus_{\lambda} (\mathfrak{b}_{\lambda} \oplus \mathfrak{c}_{\lambda})\)，其中，对每个 \(\lambda, \mathfrak{b}_{\lambda}\) 与 \(\mathfrak{c}_{\lambda}\) 都是 A 的理想。利用习题 19，有 \(\mathbf{P} = \mathbf{L} \oplus \mathbf{Q}\)，其中 \(\mathbf{L}\) 同构于 \(\mathbf{A}^{(I)}\)（I 为无限集），而 \(\mathbf{Q} = \bigoplus_{\alpha \in I} \mathfrak{a}_{\alpha}\)，其中诸 \(\mathfrak{a}_{\alpha}\) 是 A 的理想。然后应用《代数》第 II 章 § 2 习题 3。）
\end{enumerate}

¶ 22. 设 A 为局部环，m 为其极大理想，E 为有限生成 A-模。

\begin{enumerate}
\def\labelenumi{(\alph{enumi})}
\item
  若 \(r = \gamma(\mathrm{E})\) 是 \(\mathbf{E}\) 的一个生成系的最小基数，则 \(r\) 也是最小整数 \(h\)，使行列式理想 \(\mathfrak{d}_h(\mathbf{E})\) 等于 \(\mathbf{A}\)，同时也是最小整数 \(k\)，使 \(\bigwedge^{k+1} \mathbf{E} = 0\)（注意 \(r\) 是 \(\mathbf{E}/\mathfrak{m}\mathbf{E}\) 在 \(\mathbf{A}/\mathfrak{m}\) 上的秩）。
\item
  若记 \(e(E) = \mathfrak{d}_{r-1}(E)\)，证明 \(\bigwedge^{r} E\) 同构于 \(A / e(E)\)。理想 \(a\)（在 \(A\) 中）包含 \(e(E)\) 的充分必要条件是 \(E / aE\) 是自由 \((A / a)\) -模（归结为 \(a = 0\) 的情形）。
\item
  若 \(\mathfrak{e}(\mathbf{E})\) 是主理想 \(A\alpha\)，证明 E 包含一个同构于 \(A/A\alpha\) 的直和项（把 E 写成 \(A^{r}\) 对某个子模 R 的商，并注意 R 中有一个形如 \(\alpha y\) 的元素，其中 y 是 \(A^{r}\) 的某组基中的元素）。
\item
  设 \(\mathfrak{b}_h^{\prime}(\mathbf{E})\) 为 \(\bigwedge^{h+1}\mathbf{E}\) 的零化子（对 \(0 \leqslant h \leqslant r - 1\)）。证明下列条件等价：
\end{enumerate}

\((\alpha)\) 诸理想 \(\mathfrak{d}_h(\mathbf{E})\)（\(0 \leqslant h \leqslant r - 1\)）是主理想。

(β) 诸理想 \(\mathfrak{d}_h^\prime (\mathrm{E})\) \((0\leqslant h\leqslant r - 1)\) 是主理想。

\((\gamma)\) E 是 r 个同构于 A/A \(\lambda_{h}\) 的模的直和，其中 \(\lambda_{h+1}\) 整除 \(\lambda_{h}\)（对 \(0 \leqslant h \leqslant r - 1\)）。

而且，若这些条件成立，则 \(\mathfrak{d}_h'(\mathrm{E}) = \mathrm{A}\lambda_h\)（对 \(0 \leqslant h \leqslant r - 1\)），且 \(\mathfrak{d}_h(\mathrm{E}) = \lambda_h\lambda_{h+1} \ldots \lambda_{r-1}\mathrm{A}\)（对 \(0 \leqslant h \leqslant r - 1\)）。（对 \(r\) 作归纳，并利用 (c)。）

\begin{enumerate}
\def\labelenumi{(\alph{enumi})}
\setcounter{enumi}{4}
\tightlist
\item
  再设 A 是整环。证明：若 p 是 A 的素理想，使 E/pE 是自由 (A/p)-模且 \(E_{p}\) 是自由 \(A_{p}\) -模，则 E 是自由 A-模（利用 (b) 证明 \(\gamma(E_{p}) = \gamma(E)\) 且 \(e(E_{p}) = e(E)\)）。
\end{enumerate}

\begin{enumerate}
\def\labelenumi{\arabic{enumi}.}
\setcounter{enumi}{22}
\item
  设 A 为环，E 为有限生成 A-模，r 为使 \(\bigwedge^{k+1}E=0\) 的最小整数 k，\(e(E)\) 为 \(\bigwedge^{r}E\) 的零化子。证明对 A 中每个理想 \(a\supset e(E)\)，\((A/a)\) -模 E/aE 是平坦的（归结为 A 是局部环的情形并利用习题 22 (b)）。由此推出：若 r 是 A 的 Jacobson 根，且对 A 的每个极大理想 m，向量 \((A/m)\) -空间 E/mE 的秩都相同，则 E/rE 是平坦 \((A/r)\) -模。
\item
  设 A 为整闭 Noether 整环。格 M（关于 A 的）是自反的，其充分必要条件是它满足下列条件：对 A 中元素的每个有序对 \((a, b)\)，若 \(a \neq 0\) 且 A/aA 上比率为 b 的位似是单射，则 M/aM 上比率为 b 的位似也是单射。（为看条件必要，考虑 M 中满足 \(ax + by = 0\) 的两个元素 x、y，并证明对所有 \(p \in P(A)\) 有 \(y \in aM_{p}\)。为证条件充分，利用定理 2 的判别法 (c)，并注意（利用 § 1 第 4 节命题 8）假设可以写成
\end{enumerate}

\[
\operatorname{Ass} (\mathbf {M} / a \mathbf {M}) = \operatorname{Ass} (\mathbf {A} / a \mathbf {A})
\]

对 \(a \neq 0\)（\(\mathbf{A}\) 中）成立，且对每个模 \(\mathbf{E}\)（满足 \(\mathbf{M} \subset \mathbf{E} \subset \mathbf{V}\)），存在 \(a \neq 0\)（在 \(\mathbf{A}\) 中）使 \(a\mathbf{M} \subset a\mathbf{E} \subset \mathbf{M}\)。

\begin{enumerate}
\def\labelenumi{\arabic{enumi}.}
\setcounter{enumi}{24}
\tightlist
\item
  设 \(A = \bigoplus_{n \geq 0} A_n\) 为正次数的 Noether 分次环。证明：若 A 是因子分解整环，则 \(A_0\) 是因子分解整环，且诸 \(A_n\) 是自反的。（为证 \(A_0\) 是因子分解整环，利用 § 3 第 2 节定理 1 的判别法 (c)；为证诸 \(A_n\) 自反，利用习题 24。）
\end{enumerate}

¶ 26. 设 A 为 Noether 环，M 为有限生成 A-模。对称代数 S(M) 是因子分解整环的充分必要条件是：A 是因子分解整环且诸 \(S^{n}(M)\) 是自反 A-模。（为证条件充分，首先注意，若 T = A - \{0\}，则 S(M) 等同于 \(T^{-1}S(M)\) 的一个子环；然后证明 A 的每个极端元 p 在 S(M) 中也是极端元，这可归结为证明：若 p 整除 S(M) 中两个齐次元素之积 xy，则它整除二者之一；最后应用 § 3 第 4 节命题 3。）

\specialchapter{历史注记（第 I 至 VII 章）}

（方括号中的数字指本注记末尾的参考书目。）

“抽象”交换代数是一项晚近的创造，但其发展只有放在代数数论与代数几何——交换代数正是由它们孕育而生——的发展之中才能理解。

人们曾经相当有道理地猜测，Fermat 声称自己拥有的那个著名“证明”，即方程 \(x^{p} + y^{p} = z^{p}\)（p 为奇素数且 x、y、z 为整数 \(\neq 0\)）不可能成立，依赖于分解式

\[
(x + y) (x + \zeta y) \dots (x + \zeta^ {p - 1} y) = z ^ {p}
\]

在环 \(Z[\zeta]\)（其中 \(\zeta \neq 1\) 是 \(p\) 次单位根）中的分解，以及该环中（假定它是主理想整环）的一个整除性论证。无论如何，在 Lagrange 的著作中已找到一个类似论证的概要（{[}2{]}, 卷 II, 页 531）；正是通过这类论证（带有某些变化，特别是旨在降低方程次数的变数替换），Euler（{[}1{]}, 卷 I, 页 488）(*) 与 Gauss（{[}3{]}, 卷 II, 页 387）对 \(p = 3\) 证明了 Fermat 定理，Gauss（同前引）与 Dirichlet（{[}4{]}, 卷 I, 页 42）对 \(p = 5\)，而 Dirichlet 证明了方程 \(x^{14} + y^{14} = z^{14}\) 的不可能性（{[}4{]}, 卷 I, 页 190）。最后，在其关于数论的最初研究中，Kummer 曾相信他以此得到了一个一般的证明，而无疑正是这个错误（Dirichlet 曾向他指出）把他引向分圆域算术的研究，并最终从中学成功地导出他对素数 \(p < 100\) 的证明的正确版本{[}7d{]}。

另一方面，Gauss 1831 年关于双二次剩余的著名论文（其结果来自对“Gauss 整数”环 \(\mathbf{Z}[i]\) 中整除性的细致研究（{[}3{]}, 卷 II, 页 109）），清楚地显示了把整除性推广到代数数对数论经典问题可能具有的意义(*)；因此毫不奇怪，在 1830 至 1850 年间，这一理论成为德国数学家大量工作的对象，先是 Jacobi、Dirichlet 与 Eisenstein，稍后是 Kummer 及其学生兼友人 Kronecker。我们在这里不谈单位理论，它是数论中过于专门的分支，那里进展非常迅速：Eisenstein 得到了三次域单位群的结构，Kronecker 得到了分圆域单位群的结构，就在 Dirichlet 于 1846 年（{[}4{]}, 卷 I, 页 640）证明一般定理之前不久，而 Hermite 几乎已独立地达到该定理（{[}8{]}, 卷 I, 页 159）。（对整个理论居于中心地位的）素因子分解问题则显得困难得多。既然 Lagrange 已给出形如 \(x^{2} + \mathrm{D}y^{2}(x,y,\mathrm{D}\) 取整数) 的数的一些例子，其因子不具有 \(m^2 +\mathrm{D}n^2\) 的形式（{[}2{]}, 卷 II, 页 465），人们实际上已经知道，一般不能期望环 \(\mathbf{Z}[\sqrt{-\mathrm{D}} ]\) 是主理想整环，而在 Euler 的莽撞之后接踵而来的是相当的审慎；例如当 Dirichlet 证明关系式 \(p^2 -5q^2 = r^5\)（\(p,q,r\) 为整数）等价于

\[
p + q \sqrt {5} = (x + y \sqrt {5}) ^ {5}
\]

（x、y 为整数）时，他只限于指出“对许多其他素数 {[}除 5 以外{]} 有类似的定理”（{[}4{]}, 卷 I, 页 31）。在 Gauss 1831 年的论文以及 Eisenstein 关于立方剩余的工作 {[}6a{]} 中，确实有关于主理想整环 Z{[}i{]} 与 Z{[}ρ{]}（ρ = (-1 + i√3)/2，一个立方单位根）中算术的高深研究，与有理整数理论完全类似，而且至少在这些例子中，二次域中的算术与 Gauss 所发展的二元二次型理论之间的紧密联系十分明显；但对一般情形，还缺少一本“词典”，使二次域能通过对 Gauss 理论作简单翻译来处理（†）。

事实上，问题最先得到解决不是对二次域，而是对分圆域（其原因要到很久以后才清楚（参见页 585））。从 1837 年起，原本是分析学家的 Kummer 转向分圆域的算术，此后近 25 年间几乎只从事于它。与前辈们一样，他研究环 Z{[}ζ{]} 中的整除性，其中 ζ 是 p 次单位根，≠1（p 为奇素数）；他很快看出，这里同样遇到非主理想整环的环，它们阻碍了算术定律推广方面的一切进展{[}7a{]}，直到 1845 年，经过 8 年的努力，才豁然开朗，这有赖于他对“理想数”的定义（{[}7c{]} 与 {[}7d{]}）。

用现代的语言说，Kummer 所做的恰好就是定义域 \(\mathbf{Q}(\zeta)\) 上的诸赋值：它们与他的“理想素数”一一对应；这样一个因子出现在数 \(x \in \mathbf{Z}[\zeta]\) 的“分解”式中的“指数”，恰是相应赋值在 \(x\) 处的值。由于 \(x\) 的共轭元也属于 \(\mathbf{Z}[\zeta]\)，其乘积 \(\mathrm{N}(x)\)（\(x(*)\) 的“范数”）是有理整数，待定义的“理想素因子”也必须是有理素数的“因子”，而要定义它们，只需说明素数 \(q \in \mathbf{Z}\) 的“理想素因子”是什么即可。对 \(q = p\)，Kummer 实际上已经证明{[}7a{]}，主理想 \((1 - \zeta)\) 是素的，且它的 \((p - 1)\) 次幂是主理想 \((p)\)；因此这一情形不产生新问题。对 \(q \neq p\)，看来指导 Kummer 的想法是把分圆方程 \(\Phi_p(z) = 0\) 换成同余式 \(\Phi_p(u) \equiv 0 \pmod{q}\)，换句话说，把分圆多项式 \(\Phi_p(\mathbf{X})\) 在域 \(\mathbf{F}_q\) 上分解，并把这个多项式的每个不可约因子与一个“理想素因子”相对应。一个简单的情形（Kummer 在不加证明地宣布其结果的短文 {[}7b{]} 中明确提到）是 \(q \equiv 1 \pmod{p}\) 的情形；若 \(q = mp + 1\) 且 \(\gamma \in \mathbf{F}_q\) 是 1 的本原 \((q - 1)\) 次根，则在 \(\mathbf{F}_q[\mathbf{X}]\) 中

\[
\Phi_ {p} (\mathbf {X}) = \prod_ {k = 1} ^ {p - 1} \left(\mathbf {X} - \gamma^ {k m}\right)
\]

因为 \(\gamma^{pm} = 1\)。于是把每个因子 \(\mathbf{X} - \gamma^{km}\) 与 \(\mathfrak{q}_k\)（\(q\) 的一个“理想素因子”）相对应，Kummer 说：元素 \(x \in \mathbf{Z}[\zeta]\)（设 P 是它在 \(\mathbf{Q}\) 上的极小多项式）被 \(\mathfrak{q}_k\) 整除，如果在 \(\mathbf{F}_q\) 中 \(\mathrm{P}(\gamma^{km}) = 0\)；总之，用现代的语言说，他把商环 \(\mathbf{Z}[\zeta]/q\mathbf{Z}[\zeta]\) 写成了同构于 \(\mathbf{F}_q\) 的域的直和。对 \(q \not\equiv 1 \pmod{p}\)，\(\Phi_p(\mathbf{X})\) 在 \(\mathbf{F}_q[\mathbf{X}]\) 中的不可约因子不再是一次的，因此需要把 \(\mathrm{P}(\mathbf{X})\) 中的 X 换成 \(\Phi_p\) 在 \(\mathbf{F}_q[\mathbf{X}]\) 中各因子的“Galois 虚”根。Kummer 通过——按我们今天的说法——过渡到 \(q\) 的分解域 K 来避开这一困难；若 \(f\) 是使 \(q^f \equiv 1 \pmod{p}\) 的最小整数且 \(p - 1 = ef\)，则 K 恰是 \(\mathbf{Q}(\zeta)\) 中由阶为 \(f\) 的子群（属于 \(p - 1\) 阶循环 Galois 群，即 \(\mathbf{Q}(\zeta)\) 在 \(\mathbf{Q}\) 上的 Galois 群）的不变量组成的子域；换句话说，它是 \(\mathbf{Q}(\zeta)\) 中次数为 \(e\)（在 \(\mathbf{Q}\) 上）的唯一子域；自 Gauss 的《Disquisitiones》以来人们对它已很熟悉，它由“周期”

\[
\eta_ {k} = \zeta_ {k} + \zeta_ {k + f} + \zeta_ {k + 2 f} + \dots + \zeta_ {k + (e - 1) f}
\]

\((0 \leqslant k \leqslant e - 1, \zeta_{\nu} = \zeta^{g^{\nu}} \text{ 其中 } g \text{ 是同余式的一个本原根 } z^{p-1} \equiv 1 (\text{mod. } p))\) 生成，这些周期构成它的一个正规基。若 \(R(X)\) 是这些“周期” \(\eta\) 中任何一个的极小多项式（首一且具有有理整系数），Kummer 从 Gauss 的公式出发，证明在域 \(F_q\) 上 \(R(X)\) 也分解为互异的一次因子 \(X - u_j\)（\(1 \leqslant j \leqslant e\)），而他现在把每个 \(u_j\) 与一个“理想素因子” \(q_j\) 相对应。为定义“被 \(q_j\) 整除”，Kummer 把每个 \(x \in Z[\zeta]\) 写成

形式 \(x = \sum_{k=0}^{\infty} \zeta^k y_k\)，其中每个 \(y_k \in \mathbf{K}\) 本身都可唯一地写成次数 \(\leqslant e - 1\) 的 \(\eta\) 的多项式，系数为有理整数；他说 \(x\) 被 \(q_j\) 整除当且仅当用 \(u_j\) 代替 \(\eta\)（在每个 \(y_k\) 中）后，所得到的 \(\mathbf{F}_q\) 中元素全为零。但还必须定义 \(q_j\) 在 \(x\) 中的“指数”。为此，Kummer 引入了我们今天会称之为 \(q_j\) 的单值化元的东西，即一个元素 \(\rho_j \in \mathbf{K}\)，满足 \(N(\rho_j) \equiv 0 \pmod{q}\)、\(N(\rho_j) \not\equiv 0 \pmod{q^2}\)，最后还满足 \(\rho_j\) 被 \(q_j\) 整除（按上面定义的意义），但不被任何其他理想因子 \(\neq q_j\)（\(q\) 的）整除。这样一个 \(\rho_j\) 的存在性实际上已在前一年被 Kronecker 在其学位论文中证明（{[}9a{]}, 页 23）；然后记 \(\rho_j' = N(\rho_j)/\rho_j\)，Kummer 说，\(q_j\) 在 \(x\) 中的指数等于 \(h\)，如果 \(x\rho_j'^h \equiv 0 \pmod{q^h}\) 但 \(x\rho_j'^{h+1} \not\equiv 0 \pmod{q^{h+1}}\)；他当然先证明关系 \(x\rho_j' \equiv 0 \pmod{q}\) 等价于 \(x\) 被 \(q_j\) 整除（按上述意义）。这些定义一经作出，把关于“理想数”的通常整除性规律推广到 \(Z[\zeta]\) 便不再有严重的困难；而且从他的第一篇论文 {[}7c{]} 起，Kummer 甚至能够利用 Dirichlet 的“抽屉法”证明“类”与“理想因子”都只有有限多个(*)。

我们不追述 Kummer 后来关于分圆域的工作的历史，那些工作涉及类数的确定以及在某些情形对 Fermat 定理证明的应用。只提一下他在 1859 年推广其方法以（至少部分地）得到“Kummer 域” \(\mathbf{Q}(\zeta, \mu)\) 中“理想素数”的方式，其中 \(\mu\) 是不可约多项式 \(\mathrm{P}(\mathbf{X}) = \mathbf{X}^p - \alpha\) 的一个根，\(\alpha \in \mathbf{Z}[\zeta]\){[}7e{]}。有意思的是，Kummer 恰是通过把 \(\mathbf{Q}(\zeta, \mu)\) 视为取 \(\mathbf{Q}(\zeta)\) 为“基域”的循环扩张 ( \(\dagger\) ) 来考虑问题的：他从一个“理想素数” \(q\)（\(\mathbf{Z}[\zeta]\) 的）出发，假设它既不整除 \(p\) 也不整除 \(\alpha\)，而这次他（用现代术语）考察多项式 \(\overline{\mathrm{P}}(\mathbf{X}) = \mathbf{X}^p - \bar{\alpha}\) 在剩余域 \(k\) 中（即 \(\mathbf{Q}(\zeta)\) 上对应于 \(q\) 的赋值的剩余域）的性态（\(\bar{\alpha}\) 是 \(\alpha\) 在 \(k\) 中的典范像）。由于 \(\mathbf{Q}(\zeta)\) 是 \(p\) 次单位根的域，\(\overline{\mathrm{P}}\) 或在 \(k\) 上不可约，或是一次因子的乘积。在第一种情形，Kummer 说 \(q\) 在 \(\mathbf{Z}[\zeta, \mu]\) 中保持素性；在第二种情形，他引入元素 \(w_i\)（\(1 \leqslant i \leqslant p\)，属于 \(\mathbf{Z}[\zeta]\)），它们在 \(k\) 中的像是 \(\overline{\mathrm{P}}\) 的根，并把每个指标 \(i\) 与一个理想素因子 \(\mathfrak{r}_i\)（\(q\) 的）相对应；然后记 \(\mathrm{W}_i(\mathbf{X}) = \prod_{j \neq i} (\mathbf{X} - w_j)\)，他说，对多项式 \(f\)（系数在 \(\mathbf{Z}[\zeta], f(\mu)\) 中），它包含因子 \(\mathfrak{r}_i m\) 次如果

\[
f (w _ {i}) \mathrm{W} _ {i} ^ {m} (w _ {i}) \equiv 0 \quad (\mathrm{mod.} q ^ {m})
\]

但

\[
f (w _ {i}) \mathrm{W} _ {i} ^ {m + 1} (w _ {i}) \not \equiv 0 \pmod {q ^ {m + 1}}.
\]

总之，他以这种方式得到了 \(\mathbf{Q}(\zeta, \mu)\) 上的、在 \(\mathbf{Q}\) 上非分歧的那些赋值，这对他所考虑的应用已经足够。

\[
\begin{array}{c c c} \ast & \ast & \ast \\ \hline \end{array}
\]

Kummer 在最初研究其 Fermat 定理工作所引向的那些特殊域时，有幸遇到许多使这些域的研究容易得多的偶然情况。把 Kummer 的结果推广到一般情形则有相当大的困难，并且要花费多年的努力。

以 Kronecker 与 Dedekind 为主要角色，Kummer 的发现之后的 40 年间，代数数论的历史与 180 年前 Newton 与 Leibniz 关于发明无穷小演算的争衡并无不同（但幸运的是没有同样激烈的意气）。Kummer 在柏林的学生、后来又成为其同事的 Kronecker（如前所述，其学位论文曾是 Kummer 理论中的一个要点）对“理想数”抱有浓厚兴趣，目的是把它们用于他自己的研究；当我们看到他早在 1853 年（{[}9b{]}, 页 10）就宣布关于 \(\mathbf{Q}\) 的 Abel 扩张结构的一般定理，而且——这也许更为惊人——在随后的岁月里创立了复乘理论并发现了类域论的最初萌芽（{[}9c{]} 与 {[}9d{]}）时，不能不钦佩他惊人的洞察力。Kronecker 1857 年致 Dirichlet 的一封信（{[}9{]}, 卷 5, 页 418–421）表明他当时已拥有 Kummer 理论的一个推广，这一点 Kummer 本人在他的一篇著作（{[}7e{]}, 页 57）中也予确认，而 Kronecker 在 1860 至 1880 年间的论文中还将多次提到这一理论(*)。

但是，尽管那时德国数论学派中没有一位数学家不知道 Kronecker 的这些工作，后者看来只把其方法的原理传达给一个由朋友和学生组成的小圈子；而当他终于在 1881 年关于判别式的论文 {[}9e{]} 中、尤其在他 1882 年的宏篇“纪念文集” {[}9f{]} 中决定发表它们时，Dedekind 不禁表达了他的惊讶（{[}10{]}, 卷 III, 页 427），因为他原先根据听到的传闻，以为那些方法完全不同（{[}10{]}, 卷 III, 页 287）。此外，Kronecker 远不具备 Dedekind 那样出色的表述与清晰的天赋，因此毫不奇怪，构成代数数论框架的主要是后者早在 1871 年就已发表的方法；无论“未定元的添加”方法多么有意思，就数论而言，它在我们眼中几乎只是 Dedekind 方法的一个变体（参见第 VII 章 § 1 习题 31），而 Kronecker 的思想对交换代数历史获得全部重要性，主要是在另一个方向上，即朝向代数几何的方向，我们将在后面看到。

出于要到很久以后才能清楚看到的理由，任何一般理论尝试的第一个准备当然是对代数整数概念的澄清。这一点大约在 1845–50 年间完成，尽管难以准确确定其出现的日期；看来很可能，是一个对加法与乘法稳定的系（或者更确切地说，我们今天所说的有限秩 Z-代数）这一思想，或多或少自觉地导致了代数整数的一般定义：事实上，当把形如 Z{[}θ{]} 的 Z-代数限制为有限秩时——通过与由单位根生成的环 Z{[}ζ{]} 类比，后者在当时始终是算术学家关注的核心——这个定义是必然会碰到的。无论如何，当 Dirichlet（{[}4{]}, 卷 I, 页 640）、Hermite（{[}8{]}, 卷 I, 页 115 与 146）与 Eisenstein（{[}6c{]}, 页 236）彼此独立地引入代数整数概念时，他们似乎并不认为自己在处理一个新概念，也不认为对它作详细研究会有用处；只有 Eisenstein 实际证明（同前引）两个代数整数的和与积仍是代数整数，而且并未声称这一结果是独创的。

一个更为微妙的问题，是确定在哪些环中可以期望 Kummer 理论的推广。Kummer 在他的第一篇短文 {[}7b{]} 中毫不犹豫地断言，通过考虑环 \(Z[\sqrt{D}]\)（D 为整数），他能用他的方法重新得到 Gauss 的二元二次型理论；他从未发展这一想法，但确实看来，无论他还是 Dedekind 以前的任何人都未曾察觉，在环 \(Z[\sqrt{D}]\) 中素“理想”因子的唯一合成是不可能的（当 \(D \equiv 1 \pmod{4}\) 时）（尽管单位立方根的例子表明，自 Gauss 以来就被考虑的环 \(Z[\rho]\) 不同于 \(Z[\sqrt{-3}]\)）(*)。在 Dedekind 与 Kronecker 之前，人们研究的环始终是 \(Z[\theta]\) 型的，有时是某些特殊的 \(Z[\theta, \theta']\) 型的环（†）。至于 Kronecker，考虑一个代数扩张中全部整数的想法，最初可能是由代数函数域的研究提示给他的，因为在这个域中该环自然地作为“在无穷远处有限”的函数全体而出现；无论如何，他在 1881 年关于判别式（早在 1862 年已写就并向柏林科学院宣布）的论文中，坚持这些域中“整数”的这一刻画{[}9e{]}。Dedekind 没有说明他自己这一思想的来源，但在他 1871 年关于数域的第一篇出版物中，该域的全体整数所成之环在其理论中起着首要的作用；也是 Dedekind 通过引入导子的概念（{[}10c{]}），阐明了这样一个环与具有同一分式域的诸子环之间的关系。

但这并非唯一的困难。要推广 Kummer 的思想，首先必须摆脱经由分解域的路径，因为在非 Abel 域的情形它自然不可能有类似物。这个迂回乍看起来也很令人惊讶并显得矫揉造作，因为从不可约多项式 \(\Phi_p(\mathbf{X})\)（\(\mathbf{Z}[\mathbf{X}]\) 中的）出发，人们会问，是什么妨碍 Kummer 把他的思想推进到逻辑的终点，并使用当时众所周知的“Galois 虚”数理论。障碍在 Dedekind 的学生 Selling 早在 1865 年作出的一次不幸的推广尝试中更清楚地显露出来：给定不可约多项式 \(\mathbf{P} \in \mathbf{Z}[\mathbf{X}]\)，Selling 把相应的多项式 \(\overline{\mathbf{P}}(\mathbf{X})\) 在 \(\mathbf{F}_q[\mathbf{X}]\) 中分解为不可约因子；这些因子的根因此属于有限扩张 \(\mathbf{F}_r\)（\(\mathbf{F}_q\) 的）；但 Selling 为了按 Kummer 的方式定义 \(q\) 的一个“理想素因子”在 \(\mathrm{P}(\mathrm{X})\) 的分裂域中某整数里的指数，毫不犹豫地在域 \(\mathbf{F}_r\) 中谈论模 \(q\) 的幂的同余式（{[}11{]}, 页 26）；而稍后当他试图处理分歧问题时，他又对 \(\mathbf{F}_r\) “添加”形如 \(x^h = q\) 的方程的“虚根”（{[}11{]}, 页 34）。显然，这些大胆的步骤（若把有限域 \(\mathbf{F}_q\) 换成 \(q\) -进域，它们本可得到辩护）在当时只能以胡说告终。所幸，Dedekind 于 1857 年（{[}10a{]}）以“高阶同余理论”为名，以另一种形式重拾了有限域的理论(*)：他把有限域的元素解释为 \(\mathbf{Z}[\mathbf{X}]\) 中多项式关于一个“双重模”的“剩余”，该双重模由系数在 \(\mathbf{Z}[\mathbf{X}]\) 中的线性组合构成，被组合的是一个素数 \(p\) 与一个首一不可约多项式 \(\mathrm{P} \in \mathbf{Z}[\mathrm{X}]\)（这无疑对他、也和 Kronecker 一样，是他们稍后各自独立达到的模的一般观念的起源）。据他本人的自述（{[}10d{]}, 页 218），Dedekind 当初处理 \(p\) 的“理想因子”在域 \(\mathbf{Q}(\xi)\)（\(\mathrm{P} \in \mathbf{Z}[\mathrm{X}]\) 是 \(\xi\) 的极小多项式）中的问题时，似乎是如下进行的（至少在“非分歧”情形，即当 \(\overline{\mathbf{P}}\)（在 \(\mathbf{F}_p[\mathbf{X}]\) 中、与 \(\mathrm{P}\) 相应的多项式）无重根时，肯定如此）：他在 \(\mathbf{Z}[\mathbf{X}]\) 中写出

\[
\mathbf {P} = \mathbf {P} _ {1} \mathbf {P} _ {2} \dots \mathbf {P} _ {h} + p. \mathbf {G}
\]

其中诸 \(\overline{\mathbf{P}}_i\) 在 \(\mathbf{F}_p[\mathbf{X}]\) 中不可约且互异；可以假设 \(\mathbf{G}\)（在 \(\mathbf{Z}[\mathbf{X}]\) 中）不被任何 \(\mathbf{P}_i\) 整除，并且他对所有 \(i\) 记 \(\mathbf{W}_i = \prod_{j\neq i}\mathbf{P}_j\)；于是，若 \(f\in \mathbf{Z}[\mathbf{X}]\)，则称 \(f(\xi)\) 包含“理想因子” \(\mathfrak{p}_i\)（\(p\) 的）

（相应于 \(\mathbf{P}_i\)）\(k\) 次如果

\[
f \mathrm{W} _ {i} ^ {k} \equiv 0 \quad (\text { modd. } p ^ {k}, \mathrm{P})
\]

且

\[
f \mathrm{W} _ {i} ^ {k + 1} \not \equiv 0 \quad (\text { modd. } p ^ {k + 1}, \mathrm{P}).
\]

这里与 Kummer 对“Kummer 域”所循方法的关系是一目了然的，并且 Kummer 对分圆域的原始定义可以由此轻易地重新得到（例如见 Zolotareff 的工作{[}14{]}，他起初独立于 Dedekind，稍后发展了这些想法）。

然而，Dedekind 与看来也作过类似尝试的 Kronecker 都未能在此方向上更进一步，两人都受阻于分歧性所带来的困难（{[}10d{]}, 页 218 与 {[}9f{]}, 页 325）(*)。若所考虑数域 K 的整数环 A 在（Z 上）拥有一个由同一整数 \(\theta\) 的幂组成的基，则把上述方法推广到 \(\mathbf{Z}[\theta]\) 中的分歧素数并无困难（如 Zolotareff 所指出的（同前引））。但存在这样的域 K，其环 A 中没有这种类型的基；而且 Dedekind 最终甚至发现，存在这样的情形：某些素数 \(p\)（域 K 的“判别式的非常规因子”）使得，对所有 \(\theta \in \mathrm{A}\)，对 \(\theta\) 在 Q 上的极小多项式应用上述方法，会导致把多个理想因子归于 \(p\)，尽管 \(p\) 在 A 中实际上是非分歧的 ( \(\dagger\) )；他承认曾被这个出乎意料的困难长久阻滞，后来才通过从零开始创立模与理想的理论予以克服，其阐述精湛（而且已完全是现代风格，与其同代人的散漫风格形成对照），见于无疑是他杰作的、Dirichlet 数论书著名的“第十一附录”{[}10f{]}。这部作品经历了三个连续的版本，但在第一个版本（1871 年作为 Dirichlet 书第二版的“第十附录”发表{[}4 bis{]}）中方法的要点已经具备，而且几乎在一举之间，代数数论从草稿与先前的摸索过渡为一门完全成熟的学科，并且已拥有其基本工具：从一开始，数域的全体整数所成之环就被置于理论的中心；Dedekind 证明了该环在 Z 上存在一组基，并由此导出域的判别式的定义，即整数环的一组基的元素与其共轭所构成的行列式的平方；不过他在第十一附录中只对二次域给出了分歧素数的刻画（作为判别式的素因子）（{[}10f{]}, 页 202），而他自 1871 年起就已拥有一般定理(*)。这部作品的中心结果是理想分解为素因子的存在性与唯一性定理，为此 Dedekind 首先发展了一个初等的“模”理论；事实上，在第十一附录中他把这个名称保留给数域的 Z-子模，但他对之所形成的概念以及他所证明的结果，其阐述方式已可直接应用于一般的模（†）；除其他内容外，应当指出早在 1871 年就引入了“传递子”概念，它在唯一分解定理的第一个证明中（与“升链条件”一起）起了重要作用。在其后的两个版本中，Dedekind 又给出了该定理的另外两个证明，他理所当然地把这个定理视为其理论的基石。这里应当指出，正是在第三个证明中使用了分式理想（早在 1859 年 Kummer 就对分圆域引入了它们），并确立了它们成群这一事实；我们将在后面回到第二个证明（页 594）。

所有这些结果（除去语言）无疑在 1860 年前后已为 Kronecker 所知，作为他更一般概念的特例，那些概念我们将在后面谈到（而 Dedekind 承认他只是在 1869–70 年才克服了其理论的最后困难（{[}10e{]}, 页 351））( \(\ddagger\) )；就数域而言，特别应当强调，Kronecker 在当时已经知道，整个理论可以从一个“基域” \(k\)（本身是数域，异于 \(\mathbf{Q}\)）出发而无本质改变地适用，复乘理论自然导致了这一观点；他由此对某些域 \(k\) 认识到，存在代数扩张 \(\mathbf{K} \neq k\)，它们在 \(k\) 上非分歧（{[}9f{]}, 页 269），这一事实对 \(k = \mathbf{Q}\) 不可能成立（由 Hermite 与 Minkowski 对判别式的下界估计可得）。Dedekind 从未发展后一观点（尽管他在 1882 年关于差积的论文中指出了其可能性），而“相对域”理论的第一个系统阐述应归功于 Hilbert {[}16d{]}。

最后，在 1882 年（{[}10e{]}），Dedekind 通过引入差积完成了这一理论，它给了他判别式的一个新定义，并使他能够精确地确定判别式分解中诸理想素因子的指数。也是在大约这个时候，他开始对 Galois 扩张呈现的特殊性质发生兴趣，引入了分解群与惯性群的概念（在他 1894 年才发表的论文 {[}10g{]} 中），甚至（在他生前未发表的论文中（{[}10{]}, 卷 II, 页 410–411））给出了分歧群的一个概要，Hilbert（独立于 Dedekind）在稍后发展了它（{[}16c{]} 与 {[}16d{]}）。

这样，大约在 1895 年，代数数论完成了其发展的第一阶段；这一形成时期锻造的工具将使它几乎立即进入下一阶段，即一般类域论（或者等价地说，数域的 Abel 扩张理论），它延续至今，我们不在此描述。从交换代数的观点看，可以说与此同时 Dedekind 整环的历史已实际完成（撇开其公理化刻画不谈），以及这些整环上有限生成模的结构（在数域的情形，它要到 1912 年才由 Steinitz 实质性地阐明{[}20b{]}）(*)。

\begin{center}* * *\end{center}

交换代数后来的进展主要源自相当不同的问题，它们来自代数几何（在目前时期的“抽象”发展之前，它还直接影响了数论）。

我们在此不涉及代数几何的详细历史，在 Riemann 去世之前它几乎不触及我们的主题。只需回顾，它主要研究复射影平面中的代数曲线，通常以射影几何方法（用或不用坐标）处理。与之平行，Abel、Jacobi、Weierstrass 与 Riemann 发展了单复变“代数函数”及其积分的理论；数学家们显然意识到这一理论与平面代数曲线几何之间的联系，甚至有时知道可以“把分析应用于几何”；但研究代数函数所用的方法主要具有“超越”性质，在 Riemann 之前即已如此（†）；这一特征在 Riemann 的工作中更为突出，那里引入了“Riemann 面”以及在这样的曲面上定义的任意解析函数。几乎在 Riemann 去世之后立即，Roch 尤其是Clebsch 认识到，有可能从 Riemann 以超越方法得到的深刻结果获得对曲线射影几何的大量惊人的应用，这当然会促使当时的几何学家给出这些结果的纯“几何”证明；这一纲领由 Clebsch 与 Gordan 执行得并不完全，几年后由 Brill 与 M. Noether 借助对给定曲线上变动点组以及过这些点组的辅助曲线（“伴随曲线”）的研究而完成{[}13{]}。但即使对其同代人而言，Riemann 的超越方法（特别是他对拓扑概念与“Dirichlet 原理”的使用）看来也建立在不稳固的基础之上；而且尽管 Brill 与 Noether 比当时大多数“综合”几何学家更审慎（见下文页 593），他们的几何-解析方法并非无可指摘。正是为了给平面代数曲线理论一个坚实的基础，Dedekind 与 Weber 于 1882 年发表了他们关于这一主题的宏篇论文 {[}10 bis{]}：“下面发表的研究”，他们说，“旨在以同时简单、严格而且完全一般的方式，为单变量代数函数理论——Riemann 的主要创造之一——奠定基础。在此前关于这一主题的研究中，一般都对所考虑函数的奇异性作了限制性假设，而所谓的例外情形，或者作为极限情形顺带提及，或者被完全忽略。同样，若干关于连续性或解析性的基本定理被接受下来，其‘自明性’依赖于各种各样的几何直观”（{[}10 bis{]}, 页 181）(*)。

他们工作的本质思想，是把单变量代数函数理论仿照 Dedekind 刚刚发展起来的代数数论来塑造；为此，他们必须首先从“仿射”观点看问题（与其同代人相反，后者一成不变地把代数曲线视为嵌入复射影空间之中）：因此，他们从有理函数域 \(\mathbf{C}(\mathbf{X})\) 的有限代数扩张 K 以及 K 中“整代数函数”的环 A（即该域中在多项式环 C{[}X{]} 上为整的元素）出发；他们的基本结果——不使用任何拓扑考虑而得到(*)——是 A 为一个 Dedekind 整环，“第十一附录”的全部结果都可以照搬地（甚至如 Dedekind 与 Weber 所注记的、以更简单的方式应用于它，尽管他们尚未看清原因（{[}10{]}, 卷 I, 页 268））。做完这些之后，他们证明其定理实际上是双有理不变的（换句话说，只依赖于域 K），特别是不依赖于开头所作“无穷远直线”的选取。对我们来说无疑还有意思得多的是：为了定义与 K 相应的“Riemann 面”的点（特别是不能对应于 A 的理想的“无穷远点”），他们被引向引入域 K 的位的概念：他们所处的境地，与 Gelfand 1940 年创立赋范代数理论时所处的境地相同——知道一个元素集 K，其元素并非一开始就作为函数给出，却又想把它们当作函数看待；而且，为得到这些假想函数的定义集，他们第一次有了这样的想法（Gelfand 承袭了它，而且它因在现代数学中随时被使用而变得司空见惯）：把集合 E 的每点 x 与从 E 到集合 G 的映射所成之集 F 相对应，得到 F 到 G 的映射 \(f \mapsto f(x)\)，换句话说，在表达式 \(f(x)\) 中把 f 当作变量而把 x 固定，这与经典传统相反。最后，他们从位的概念出发毫无困难地定义了“正除子”（他们的术语中称“多边形”），它们以 A 的理想为特例，并与 Brill 与 Noether 的“点组”相对应；但是，尽管他们把主除子与微分的除子写作正除子的“商”，他们却没有给出除子的一般定义，直到 1902 年，Hensel 与 Landsberg 才通过与分式理想类比引入这一概念，它将一直使纯“几何”方法的拥护者们为难（他们不得已以“虚系统”之名定义这些除子，却又因不能给它们“具体”解释而局促不安）。

同年 1882 年还出现了人们等待了 20 多年的 Kronecker 的宏篇论文 {[}9f{]}。它比 Dedekind 与

Weber 的工作远为宏大，但不幸也远为含糊与晦暗。它的中心主题是（用现代语言说）研究多项式环 \(C[X_{1},\ldots,X_{n}]\) 或 \(Z[X_{1},\ldots,X_{n}]\) 之一上的有限整代数的理想；Kronecker 先验地限于有限生成的理想（它们全为有限生成这一事实，只是几年后才由 Hilbert 在其不变量理论工作中（对 \(C[X_{1},\ldots,X_{n}]\) 的理想）证明{[}16a{]}）。就 \(C[X_{1},\ldots,X_{n}]\) 或 \(Z[X_{1},\ldots,X_{n}]\) 而言，这自然导致把这两个环中每个理想与由该理想全部元素公共零点组成的“代数簇”相联；而 19 世纪对 2 维与 3 维几何的研究会直觉地引出这样的想法：每个簇都是有限多个“不可约”簇的并，这些簇的“维数”未必全都相同。看来这一事实的证明正是 Kronecker 为自己设定的目标，尽管他从未明言，其论文中既找不到“不可约簇”的定义，也找不到“维数”的定义。事实上，他只约略指出，一个一般的消去法(*)如何从所考虑理想的生成系出发给出有限多个代数簇，对其中每个簇，在适当的坐标系中若干坐标是任意的，其余坐标则是它们的“代数函数”（†）。但若 Kronecker 确以分解为不可约簇为目标，那就必须承认，他只在主理想这一初等情形达到目的，在该情形他有效地证明——推广了 Z{[}X{]} 上 Gauss 的经典引理（{[}3{]}, 卷 I, 页 34）——整环 \(C[X_{1},\ldots,X_{n}]\) 与 \(Z[X_{1},\ldots,X_{n}]\) 是因子分解的；而在一般情形，Kronecker 是否拥有素理想的概念是可疑的（他称为“Primmodulsystem”的东西是一个不能分解为两个其他理想之积的理想（{[}9f{]}, 页 336）；这尤其令人惊讶，因为 Dedekind 1871 年已给出的定义是完全一般的）。

然而必须说，Kronecker 的消去法若运用得当，确实导出代数簇分解为其不可约分支：这一点由 E. Lasker 在其 1905 年关于多项式理想的宏篇论文的开头清楚地建立{[}19{]}；他正确定义了（\(\mathbf{C}^n\) 中的）不可约簇概念：代数簇 V 使得，两个多项式的积只有在其中一个于整个簇 V 上为零时才能在整个 V 上为零，他还给出一个与坐标轴选取无关的定义。在他插入这项工作的有趣的历史性考虑中，Lasker 表明他不仅对 Kronecker 与 Dedekind 的纯粹代数倾向感兴趣，而且对 Clebsch 与 M. Noether 学派的几何方法所提出的问题感兴趣，特别是后者 1873 年证明的著名定理{[}12{]}。用今天的话说，他实质上关心的是确定理想 \(a\)，它由 \(\mathbf{C}[\mathbf{X}_1,\ldots ,\mathbf{X}_n]\) 中在 \(\mathbf{C}^n\) 的给定集合 M 的各点处为零的多项式组成；通常 M 是有限多个多项式 \(f_{i}\) 的公共零点的“代数簇”，而很长一段时间里，人们似乎默认（当然没有论证）至少对 \(n = 2\) 或 \(n = 3\)，理想 \(a\) 就由诸 \(f_{i}\) 生成(*)。M. Noether 已证明，即使对 \(n = 2\) 及两个多项式 \(f_{1},f_{2}\)，这一般也不真，他并给出了 \(a\) 由 \(f_{1}\) 与 \(f_{2}\) 生成的充分条件。十年后，Netto 证明，在对 \(f_{1}\) 与 \(f_{2}\) 不作任何假设时，\(a\) 的某个幂总包含于 \(f_{1}\) 与 \(f_{2}\) 生成的理想{[}15{]}，Hilbert 1893 年在其著名的零点定理中推广了这个定理{[}16b{]}。无疑受这一结果的启发，Lasker 在其论文中引入了准素理想的一般概念 ( \(\dagger\) )，于环 \(\mathbf{C}[\mathbf{X}_1,\dots ,\mathbf{X}_n]\) 与 \(\mathbf{Z}[\mathbf{X}_1,\dots ,\mathbf{X}_n]\) 中（在先通过转录 Dedekind 的定义给出这些环的素理想定义之后），并证明（ \(\ddagger\) ）这两个环中每个理想都存在准素分解(*)。他似乎并不关心这一分解中的唯一性问题；是 Macaulay 稍后{[}21{]}引入了“浸入”与“非浸入”准素理想的区分，并证明后者唯一确定，前者则不然。最后应当指出，Lasker 还利用 Weierstrass 的“预备定理”把其结果推广到一点邻域中的收敛幂级数环。他论文的这一部分，无疑是这个环第一次从纯代数观点得到考虑，而 Lasker 在此场合发展的方法，对 Krull 1938 年创立局部环的一般理论产生了强烈的影响（参见{[}29d{]}, 页 204 及以下各处）。

\begin{center}* * *\end{center}

孕育现代交换代数的思想运动大约在 1910 年开始成形。如果说域的一般概念在 20 世纪初已经达成，那么相反，定义了环的一般概念的第一部工作大概就是 Fraenkel 1914 年的{[}23{]}。这时，作为环的例子，已有的不仅有数论与代数几何的整环，还有（形式与收敛的）幂级数环，最后还有（交换或非交换的）基域上的代数。然而，无论对域论还是对环论，催化剂的角色看来都是由 Hensel 的 p-进数理论扮演的，Fraenkel 以及 Steinitz {[}20a{]} 都特别提到它是其研究的出发点。

Hensel 关于这一主题的第一部出版物可追溯到 1897 年；他在其中从 Dedekind 与 Weber 所指出的、代数函数域 K 的 Riemann 面的点与数域 k 的素理想之间的类比出发；他提出把“Puiseux 展开”（自 19 世纪中叶起即为经典）搬到数论中，该展开在 K 的 Riemann 面的任一点附近，使每个元素 \(x \in K\) 都可表成所考虑点处“单值化元”的幂的收敛级数（负指数项只有有限多个）。Hensel 类似地证明，若 p 是 k 中位于素数 p 之上的素理想，则对每个 \(x \in k\) 可联结一个“p-进级数”，形如 \(\sum_{i} \alpha_{i} p^{i} \left( \text{或 } \sum_{i} \alpha_{i} p^{i/e} \text{若 } p \text{ 分歧于 } p \right)\)，其中诸 \(\alpha_{i}\) 取自理想 p 的剩余域的一个给定代表系；但他的伟大独创性在于想到，即使这些“展开”不对应 k 中任何元素，也可以考虑它们，这与 Riemann 面上超越函数的整级数展开相类比{[}18a{]}。

在其生涯的其余时间里，Hensel 致力于逐步打磨和完善他的新演算；如果说他的方式在我们看来迟疑而笨重，那么不应忘记，至少在起初，他手头没有现代数学的任何拓扑或代数工具，而这些工具本可使其任务轻松。在其最初的出版物中，他很少谈及拓扑概念，总的来说，对他而言（用现代术语说）p-进整数环（p 为数域 k 的整数环 A 中的素理想）就是环 A/p \(^{n}\)（n 无限增大）的逆极限，这纯粹是在代数意义上说的；而为建立该环及其分式域的性质，每一步都不得不使用多少有些费力的特设论证（例如为证明 p-进整数成一整环）。把拓扑概念引入 p-进域的想法，1905 年之前未见于 Hensel 的著作{[}18d{]}；直到 1907 年，在出版了按其思想重述代数数论的书{[}18f{]}）之后，他才达到 p-进绝对值的定义及其本质性质{[}18e{]}，由此出发，他得以仿照 Cauchy 的理论发展一种新的“p-进分析”，并能在数论中富有成效地加以应用（特别是使用 p-进指数与对数），其重要性自此与日俱增。

Hensel 从一开始就清楚地看到其理论给经典论述带来的简化：它使问题得以“局部化”，使工作得以在一个其中整除性质成为平凡的域中进行，而且多亏他早在 1902 年发现的基本引理{[}18c{]}，对“约化的”多项式 mod p 无重根的多项式的研究，被归结为对有限域上多项式的研究。早在 1897 年{[}18b{]}他就给出了这些简化的惊人例子，特别是在与判别式有关的问题上（尤其包括对他几年前给出的“非常规因子”存在性判别法的一个简短证明）。但在很长一段时间里，p-进数似乎引起当时数学家的巨大不信任；这无疑是对“过于抽象”观念的常见态度，但也部分地为其作者相当过分的热情所证实（这种热情在数学中新理论的狂热者身上如此常见）。Hensel 不满足于把其理论富有成效地应用于代数数，他与所有同代人一样对 e 与 \(\pi\) 的超越性证明印象深刻，或许还被同时加于数与函数的“超越的”这一形容词所误导，竟至认为他的 p-进数与超越实数之间存在联系，并一度以为他得到了 e 乃至 \(e^{e}\) 的超越性的一个简单证明（{[}18d{]}, 页 556）(*)。

1910 年之后不久，局面随新一代的崛起而改变，他们受 Fréchet 与 F. Riesz 的拓扑思想以及 Steinitz 的代数思想影响，从一开始就献身于“抽象”；这一代人将懂得如何消化 Hensel 的工作并把它安置到其真正的位置。早在 1913 年，Kürschak {[}22{]} 就给出绝对值概念的一般定义，认识到超度量绝对值（p-进绝对值即其一例）的重要性，（模仿实数情形的证明）证明了域关于一个绝对值的完备化的存在性，尤其一般地证明了把绝对值扩张到给定域的任一代数扩张的可能性。但他没有看出，绝对值的超度量特征已在素域中显露；这一点由 Ostrowski 建立，域 Q 上全部绝对值的确定，以及把具非超度量绝对值的域刻画为 C 的子域的基本定理，也应归功于他{[}24{]}。在 1920 至 1935 年间，理论又通过对未必离散的绝对值的更细致研究而趋于完备，其中包括对过渡到代数或超越扩张时出现的各种情形的考察（Ostrowski、Deuring、F. K. Schmidt）；另一方面，1931 年 Krull 引入并研究了赋值的一般概念{[}29b{]}，随后几年它将被 Zariski 及其代数几何学派大量使用（†）。这里还应提到——尽管超出我们的范围——关于完备赋值域与完备局部环结构的更深入的研究，它们属于同一时期（Hasse-Schmidt、Witt、Teichmüller、I. Cohen）。

\begin{center}* * *\end{center}

上面（页 594）提到的 Fraenkel 的工作只处理了一种非常特殊的

环类型（Artin 的、只有唯一素理想而且该素理想还被假定为主理想的环）。除 Steinitz 关于域的工作 {[}20a{]} 外，关于一般交换环研究的第一批重要工作，是 E. Noether 关于理想论的两部宏篇论文：1921 年的那部{[}25a{]}讨论准素分解，以完全的一般性重拾并在许多方面完成了 Lasker 与 Macaulay 的结果；1927 年的那部公理化地刻画了 Dedekind 整环{[}25b{]}。正如 Steinitz 对域所表明的，在这些论文中可以看到，少数几个抽象观念——如不可约理想的概念、链条件以及整闭整环的思想（如前所述，后两者已由 Dedekind 揭示）——如何仅凭自身就能导出一般的结果，而在先前已知它们的情形中，这些结果似乎与纯计算结果不可分解地纠缠在一起。

E. Noether 的这些论文，加上稍后 Artin-van der Waerden 关于除性理想{[}31{]}的工作以及 Krull 把这些理想与本性赋值相联系{[}29b{]}的工作，一个世纪前开始的(*)关于理想分解的漫长研究就此完成，与此同时现代交换代数宣告诞生。

此后不可胜数的关于交换代数的研究工作，最容易按几个重要的发展方向来归类：

\textbf{(A) 局部环与拓扑}\par

尽管萌芽已含于此前关于数论与代数几何的全部工作之中，局部化的一般思想却显露得非常缓慢。分式环的一般概念直到 1926 年才由 E. Noether 的学生 H. Grell 定义，而且只对整环{[}28{]}；它向更一般环的扩张要到 1944 年才由 C. Chevalley 对 Noether 环给出，1948 年由 Uzkov 对一般情形给出。直到大约 1940 年，Krull 及其学派几乎是仅有的在一般论证中使用对局部环 \(\mathbf{A}_{\mathfrak{p}}\)（整环 \(\mathbf{A}\) 的）的考虑的人；这些环只是从 1940 年起才随 Chevalley 与 Zariski 的工作开始显式出现于代数几何（†）。

局部环本身的一般研究始于 1938 年 Krull 的宏篇论文{[}29d{]}。这项工作最重要的结果涉及维数理论与正则环，我们不在此讨论；但任意 Noether 局部环的完备化在此第一次出现，还有与局部环相伴的分次环的一个尚不完善的形式(*)；后者直到大约 1948 年才由 P. Samuel{[}36{]}以及 Leray 与 H. Cartan 在代数拓扑方面的研究独立地定义。Krull 在上述提到的工作中几乎不用拓扑语言；但早在 1928 年{[}29a{]}，他就证明在 Noether 环 A 中，一个理想 a 的诸幂之交是这样的 \(x \in A\) 所成之集：\(x(1 - a) = 0\) 对某个 \(a \in a\) 成立；由此容易推出，对 A 的每个理想 m，A 上的 m-进拓扑在理想 a 上诱导出 a 上的 m-进拓扑；在 1938 年的论文中，Krull 补全了这一结果，证明在 Noether 局部环中每个理想都是闭的。这些定理随后由 Chevalley 推广到 Noether 半局部环，再由 Zariski 推广到冠以其名的环{[}33b{]}；拓扑环中“线性紧”概念的引入，以及完备半局部环结构的确定，也应归功于 Chevalley{[}32b{]}。

\textbf{(B) 由局部到整体的过渡}\par

自 Weierstrass 以来，人们习惯把单变量解析函数（特别是代数函数）与它在 Riemann 面上其有定义的所有点处的“展开”全体相联。在其数论书的引言中（{[}18f{]}, 页 V），Hensel 类似地把代数数域 k 的每个元素与它在 k 关于其上全部绝对值的诸完备化中相应元素所成之集相联 \(k (\dagger)\)。可以说，在现代交换代数中，正是这一观点取代了理想表为素理想之积的分解公式（在某种意义上推广了 Kummer 最初的观念）。Hensel 的评述相当于隐式地把 k 嵌入其全部完备化之积；Chevalley 1936 年以其“伊代尔”理论{[}32a{]}显式地这样做了，它完善了 Prüfer 与 von Neumann 较早的类似想法（后者只限于把 \(k\) 嵌入其诸 p-进完备化之积）(*)。虽然这稍稍超出我们的范围，但在此重要的是指出，多亏伊代尔群上一个适当的拓扑，局部紧群的全部技巧（包括 Haar 测度）可以由此非常有效地应用于数论。

在更一般的脉络中，Krull 把整闭整环刻画为赋值环之交的定理{[}29b{]}（它也相当于把所考虑的整环嵌入赋值环之积）常能方便这些环的研究，尽管该方法只对 Krull 整环的本性赋值真正可行。此外，Krull 常给出{[}29e{]}“由局部到整体”方法的（十分初等的）例子，该方法通过就整环 A 的全部素理想验证某性质来对 A 的诸“局部化”环 \(A_{p}\) 证明该性质（†）；更近的时候，Serre 察觉到这一方法对任意交换环 A 都行得通，它还可应用于 A-模，而且甚至常常只需在 A 的极大理想处“局部化”（第 II 章 § 3 定理 1）：这一观点与关于“谱”以及定义在这些谱上的层的思想紧密相连（见下文页 602）。

\textbf{(C) 整数与整闭包}\par

我们看到，最先为数域引入的代数整数概念，已被 Kronecker 与 Dedekind 推广到代数函数域，尽管在此情形它可能显得相当人为（不对应于任何射影概念）。E. Noether 1927 年的论文，继之以 Krull 自 1931 年起的工作，表明了这些概念对更一般的环所呈现的意义（ \(\ddagger\) ）。特别是，把素理想提升到整代数的定理{[}29c{]}应归功于 Krull，把 Dedekind-Hilbert 的分解群与惯性群理论加以推广{[}29b{]}也应归功于他。至于 E. Noether，我们得益于她正规化引理的一般表述(*)（除其他结果外，Hilbert 零点定理即由之导出），以及一个整闭整环的整闭包在该整环上有限的首个一般判别法（转录了 Kronecker 与 Dedekind 的经典论证）。

最后，这里应当指出，整闭整环概念当今之所以重要，原因之一归功于 Zariski 关于代数簇的研究；他发现“正规”簇（即其诸局部环为整闭整环的那些簇）以特别令人愉快的性质为标志，特别是它们没有“余维数 1 的奇点”这一事实；随后人们看到，类似的现象对“解析空间”也为真。于是“正规化”（即对簇的诸局部环取其整闭包的运算）已成为现代代数几何武库中的一件强大武器。

\textbf{(D) 模的研究与同调代数的影响}\par

E. Noether 与 W. Krull 在代数中的工作的一个显著特征，是“线性化”的倾向，它延伸了 Dedekind 与 Steinitz 给予域理论的类似发展；换句话说，理想首先被当作模来考虑，于是线性代数的全部构造（商、乘积，以及更近的张量积与同态模的构成）都施于其上，一般产生的模不再是理想。由此很快看出，在许多问题中（对交换或非交换环），不应只限于研究环 A 的理想，而相反应当对 A-模一般地陈述定理（有时附加某些有限性条件）。

同调代数的介入有力地强化了上述倾向，因为代数的这一分支本质上处理线性性质的问题。我们不在此重述其历史；但有意思的是指出，同调代数的几个基本概念（如投射模的概念与 Tor 函子的概念）是在对 Dedekind 整环上的模关于张量积的性态的细致研究中诞生的，该研究由 H. Cartan 于 1948 年进行。

反过来，可以预见，同调代数作为 Ext 函子（投射模与内射模）和 Tor 函子（平坦模）的“泛零化子”而自然引入的新模类，会给交换代数投来新的光芒。事实上，主要是投射模以及更多的平坦模显示出了用处：后者的重要性尤其源于 Serre {[}38b{]} 首先作出的评注，即局部化与完备化自然地引入平坦模，从而以令人满意得多的方式“解释”了这些运算已知的性质，使它们更便于使用。此外还应提到（如我们在后面各章将看到的），同调代数的应用远不限于这些，它在代数几何中正起着越来越重要的作用。

\textbf{(E) 谱的概念}\par

日期上最新的交换代数新概念有一段复杂的历史。Hilbert 的谱定理引入了 Hilbert 空间的正交投影算子的有序集，它们构成一个“Boole 代数”（或更确切地说一个 Boole 格）(*)，与 \(\mathbf{R}\) 的（关于适当测度的）可测子集的类所成的一个 Boole 格一一对应。无疑是他约 1935 年关于 Hilbert 空间上算子的早期工作，促使 M. H. Stone 一般地研究 Boole 格，特别是寻求用一集合的子集（或关于某个等价关系的子集类）给出它们的“表示”。他观察到，若在 Boole 格上以 \(xy = \inf(x, y)\) 定义乘法、以 \(x + y = \sup(\inf(x, y'), \inf(x', y))\) 定义加法，它就成为一个交换环（而且还是非常特殊的一类环）。在所论 Boole 格恰为集合 \(\mathfrak{P}(\mathbf{X})\)（有限集 \(\mathbf{X}\) 的全部子集所成之集）这一特殊情形中，立即可见 \(\mathbf{X}\) 的元素与相应的“Boole”环的极大理想自然地一一对应；而 Stone 正是通过类似地考虑相应环的极大理想所成之集，并把 Boole 格的每个元素与包含它的极大理想所成之集相联，得到他关于 Boole 格的一般表示定理{[}30a{]}。

另一方面，拓扑空间的开且闭子集全体是 Boole 格的一个众所周知的经典例子。在第二篇论文{[}30b{]}中，Stone 证明事实上每个 Boole 格也都同构于这种类型的一个 Boole 格。为此当然必须在“Boole”环的极大理想所成之集上定义一个拓扑；这很简单地这样完成：对每个理想 a，取包含 a 的极大理想所成之集作为闭集。

我们不在此谈这些思想对泛函分析的影响，它们在 I. Gelfand 及其学派所发展的赋范代数理论的诞生中起了重要作用。但 1945 年，Jacobson 察觉到{[}34{]}，Stone 所发明的定义拓扑的过程，事实上可应用于任一环 A（交换或非交换），只要所取的理想集不是极大理想之集，而是双边“本原”理想之集（即使得 A/b 为本原环的双边理想 b）；对交换环，这些当然恰好就是极大理想。就 Zariski 而言，他在 1944 年{[}33a{]}用类似的方法在代数函数域的位所成之集上定义了一个拓扑。然而，在大多数代数学家眼中，这些拓扑长期只是奇物，原因在于它们通常不是 Hausdorff 的，而且人们对在如此异常的对象上工作有一种完全可以理解的反感。这种不信任只是当 A. Weil 于 1952 年证明每个代数簇都可赋予上述类型的一个自然拓扑、且这一拓扑允许与可微或解析流形的情形完全类比地定义纤维丛的概念{[}37{]}时，才被克服；不久之后，Serre 想到把凝聚层理论推广到这样赋予拓扑的簇，借此，在“抽象”簇的情形，拓扑起到基域为 C 时通常拓扑所起的作用，特别就应用代数拓扑的方法而言（{[}38a{]} 与 {[}38b{]}）。

从此，在整个交换代数中使用这种几何语言是自然的事。很快看出，考虑极大理想通常不足以得到有用的论断(*)，恰当的概念是把环的素理想全体之集按同样方式拓扑化。随着谱概念的引入，如今存在一部词典，使交换代数的每个定理都能用一种与 Weil-Zariski 时期的代数几何非常接近的几何语言来表达；这还立即引起了后者论域的相当大的扩展，以致交换代数几乎只是其中最初等的部分{[}39{]}。
